\setcounter{section}{4}
\setlength{\parindent}{0pt}  % niente rientro
\setlength{\parskip}{6pt}
\section{Autovalutazione}

Complessivamente, il gruppo ha lavorato in maniera solida, nonostante le criticità logistiche e la riduzione dell'organico in fase iniziale ha dimostrato una notevole resilienza e capacità di adattamento.
Entrambi i membri hanno dimostrato un atteggiamento entusiasta in ogni fase del progetto, garantendo disponibilità costante e confronto nelle fasi decisionali.
Entrando nel dettaglio delle dinamiche interne \textbf{Enrico Sartori} ha sostenuto una grossa mole di lavoro nello sviluppo applicativo e della stesura dei deliverable tecnici, dimostrando interesse nel contribuire attivamente anche nella stesura della documentazione descrittiva. Rappresentando un solido punto di riferimento organizzativo per il gruppo.
\textbf{Stefano Pezzetta} pur avendo dedicato un monte ore inferiore alla scrittura diretta del codice, ha partecipato attivamente al processo di sviluppo svolto da Enrico, richiedendo frequenti sessioni di revisione che hanno garantito la qualità e coerenza del prodotto finale rispetto alla documentazione descrittiva e alla modellazione del sistema da lui definita.

In conclusione, il flusso di lavoro si è mantenuto fluido e privo di blocchi critici.
È emersa in particolare la spiccata attitudine organizzativa di \textbf{Enrico Sartori}, che ha saputo gestire al meglio il flusso operativo, orchestrando lo sviluppo tecnico ed esprimendo opinioni chiave anche negli ambiti non propriamente di sua pertinenza.

\begin{table}[H]
\centering
\renewcommand{\arraystretch}{1.5}
\begin{tabular}{|l|c|}
\hline
\textbf{Componente del team} & \textbf{Voto proposto} \\ \hline
Enrico Sartori      & [XX] \\ \hline
Stefano Pezzetta    & [XX] \\ \hline
\end{tabular}
\caption{Proposta di valutazione individuale}
\label{tab:voti}
\end{table}